\documentclass[12pt, a4paper]{article}

\usepackage[utf8]{inputenc}
\usepackage[T1]{fontenc}
\usepackage[margin=2.5cm]{geometry}
\usepackage{graphicx}
\usepackage{hyperref}
\usepackage{listings}
\usepackage{xcolor}
\usepackage{titlesec}
\usepackage{parskip}
\usepackage{booktabs}
\usepackage{float}
\usepackage{enumitem}

% Code style
\definecolor{codebg}{RGB}{245,245,245}
\definecolor{codeframe}{RGB}{200,200,200}
\definecolor{keyword}{RGB}{0,0,180}
\definecolor{string}{RGB}{163,21,21}
\definecolor{comment}{RGB}{100,100,100}

\lstdefinestyle{code}{
    backgroundcolor=\color{codebg},
    frame=single,
    rulecolor=\color{codeframe},
    basicstyle=\ttfamily\small,
    keywordstyle=\color{keyword}\bfseries,
    stringstyle=\color{string},
    commentstyle=\color{comment}\itshape,
    breaklines=true,
    breakatwhitespace=true,
    tabsize=2,
    showstringspaces=false,
    captionpos=b,
}

\lstdefinelanguage{yaml}{
    keywords={true, false, null},
    sensitive=false,
    morestring=[b]",
    morestring=[b]',
    morecomment=[l]{\#},
}

\hypersetup{
    colorlinks=true,
    linkcolor=black,
    urlcolor=blue,
    citecolor=black
}

\titleformat{\section}{\large\bfseries}{}{0em}{}[\titlerule]
\titleformat{\subsection}{\normalsize\bfseries}{}{0em}{}

\begin{document}

% ─── Title Page ──────────────────────────────────────────────────────────────
\begin{titlepage}
    \centering
    \vspace*{2cm}
    {\Huge\bfseries Robustness \& Availability\par}
    \vspace{0.5cm}
    {\Large Show \& Tell Document\par}
    \vspace{0.5cm}
    {\large Trade2 Platform — Backend\par}
    \vfill
    {\large Moustafa Haroun\par}
    \vspace{0.3cm}
    {\large \today\par}
\end{titlepage}

\tableofcontents
\newpage

% ─── 1. Introduction ─────────────────────────────────────────────────────────
\section{Introduction}

This document describes the robustness and availability measures implemented in the Trade2 backend. The backend is a NestJS REST API deployed on a VPS using Dokploy and Docker. The measures cover four pillars:

\begin{enumerate}
    \item Zero-downtime deployments via CI/CD and Dokploy rolling updates
    \item System monitoring with Prometheus and Grafana
    \item Data consistency through database transactions
    \item Structured logging with Loki and Grafana
\end{enumerate}

% ─── 2. Zero-Downtime Deployments ────────────────────────────────────────────
\section{Zero-Downtime Deployments}

\subsection{Overview}

The application must remain available during updates. This is achieved by combining a GitHub Actions CI/CD pipeline with Dokploy, which manages rolling container updates.

\subsection{CI/CD Pipeline}

Every time a GitHub release is published, the \texttt{Build \& Publish Backend} workflow runs automatically. The workflow is defined in \texttt{.github/workflows/backend-publish.yml} and consists of two sequential jobs.

\begin{figure}[H]
    \centering
    \fbox{\rule{0pt}{5cm}\rule{12cm}{0pt}}
    \caption{GitHub Actions workflow — both jobs and their dependency}
    \label{fig:cicd}
\end{figure}

\textbf{Job 1 — build-and-push}

\begin{itemize}
    \item Checks out the repository
    \item Authenticates with the GitHub Container Registry (GHCR)
    \item Extracts image metadata: a \texttt{sha-} prefixed tag, a semver tag from the release, and \texttt{latest}
    \item Builds the Docker image using the \texttt{production} stage of the multi-stage \texttt{Dockerfile}
    \item Pushes the image to \texttt{ghcr.io/moustafaharoun/aad-backend}
    \item Uses GitHub Actions cache (\texttt{type=gha}) to speed up repeated builds
\end{itemize}

\textbf{Job 2 — deploy}

The deploy job has a \texttt{needs: build-and-push} dependency, meaning it only runs after the image is successfully pushed. It calls the Dokploy API via the \texttt{benbristow/dokploy-deploy-action} action, which triggers a re-pull and redeploy of the application container.

\begin{lstlisting}[style=code, language=yaml, caption=deploy job in backend-publish.yml]
deploy:
  runs-on: ubuntu-latest
  needs: build-and-push
  steps:
    - name: Dokploy Deployment
      uses: benbristow/dokploy-deploy-action@0.2.2
      with:
        api_token: ${{ secrets.DOKPLOY_AUTH_TOKEN }}
        application_id: ${{ secrets.DOKPLOY_APPLICATION_ID }}
        dokploy_url: ${{ secrets.DOKPLOY_URL }}
        service_type: application
\end{lstlisting}

\subsection{Rolling Updates with Dokploy}

Dokploy performs a rolling update when a redeploy is triggered. It starts a new container from the freshly pulled image and only removes the old container once the new one is healthy. This means there is no gap in availability between the old and new version.

\subsection{Database Migrations on Startup}

Schema changes are applied automatically when the new container starts. TypeORM is configured with \texttt{migrationsRun: true}, which runs all pending migration files from \texttt{dist/database/migrations/} before the application begins accepting requests. Migration files have an \texttt{up} and \texttt{down} method, making every schema change reversible.

\begin{lstlisting}[style=code, language=yaml, caption=TypeORM module configuration]
TypeOrmModule.forRoot({
  ...
  migrations: ['dist/database/migrations/*.js'],
  migrationsRun: true,
  synchronize: false,
})
\end{lstlisting}

\begin{figure}[H]
    \centering
    \fbox{\rule{0pt}{5cm}\rule{12cm}{0pt}}
    \caption{Migration log output on container startup}
    \label{fig:migrations}
\end{figure}

% ─── 3. Monitoring ───────────────────────────────────────────────────────────
\section{Monitoring}

\subsection{Overview}

The monitoring stack runs as a separate Docker Compose application in the \texttt{monitor/} directory. It consists of three services: Prometheus (metrics collection), Loki (log aggregation), and Grafana (visualisation).

\subsection{Prometheus — Metrics Collection}

Prometheus scrapes the \texttt{/metrics} endpoint of the NestJS application every 5 seconds. The application exposes metrics using the \texttt{@willsoto/nestjs-prometheus} package and \texttt{prom-client}.

\begin{lstlisting}[style=code, language=yaml, caption=prometheus.yml scrape configuration]
scrape_configs:
  - job_name: 'trade2-nestjs'
    metrics_path: '/metrics'
    scrape_interval: 5s
    static_configs:
      - targets: ['trade2-backend-k2jfhy:3000']
        labels:
          application: 'Trade2 NestJS'
\end{lstlisting}

\textbf{Collected metrics}

Two custom metrics are recorded by the \texttt{HttpMetricsInterceptor}, which wraps every incoming HTTP request:

\begin{table}[H]
    \centering
    \begin{tabular}{@{}lll@{}}
        \toprule
        \textbf{Metric} & \textbf{Type} & \textbf{Labels} \\
        \midrule
        \texttt{http\_requests\_total} & Counter & method, path, status \\
        \texttt{http\_request\_duration\_seconds} & Histogram & method, path, status \\
        \bottomrule
    \end{tabular}
    \caption{Custom Prometheus metrics}
\end{table}

In addition, Prometheus collects standard Node.js runtime metrics via \texttt{prom-client}, including heap memory usage, active handles, and event loop lag.

\subsection{Grafana Dashboard}

Grafana is pre-provisioned with a datasource configuration and a custom NestJS dashboard so it is ready immediately after deployment, without any manual setup.

\begin{table}[H]
    \centering
    \begin{tabular}{@{}llp{6cm}@{}}
        \toprule
        \textbf{Panel} & \textbf{Source} & \textbf{Query} \\
        \midrule
        Request Rate & Prometheus & \texttt{sum(rate(http\_requests\_total[1m])) by (method, path)} \\
        Error Rate 4xx / 5xx & Prometheus & \texttt{sum(rate(http\_requests\_total\{status=\textasciitilde"4.."\}[1m]))} \\
        Request Duration (p50/p95/p99) & Prometheus & \texttt{histogram\_quantile(...)} \\
        Node.js Heap Usage & Prometheus & \texttt{nodejs\_heap\_size\_used\_bytes} \\
        Active Handles & Prometheus & \texttt{nodejs\_active\_handles\_total} \\
        Event Loop Lag & Prometheus & \texttt{nodejs\_eventloop\_lag\_seconds} \\
        Total Requests (last 1h) & Prometheus & \texttt{sum(increase(http\_requests\_total[1h]))} \\
        Error \% (last 1h) & Prometheus & ratio of 4xx/5xx to total \\
        Application Logs & Loki & \texttt{\{app="trade2-nestjs"\}} \\
        Error Logs & Loki & \texttt{\{app="trade2-nestjs"\} |= `error`} \\
        \bottomrule
    \end{tabular}
    \caption{Grafana dashboard panels}
\end{table}

\begin{figure}[H]
    \centering
    \fbox{\rule{0pt}{6cm}\rule{12cm}{0pt}}
    \caption{Grafana dashboard showing request rate, error rate, and duration percentiles}
    \label{fig:grafana-metrics}
\end{figure}

\begin{figure}[H]
    \centering
    \fbox{\rule{0pt}{6cm}\rule{12cm}{0pt}}
    \caption{Grafana dashboard showing application logs and error logs panels}
    \label{fig:grafana-logs}
\end{figure}

% ─── 4. Data Consistency ─────────────────────────────────────────────────────
\section{Data Consistency}

\subsection{Database Transactions}

TypeORM's \texttt{DataSource.transaction()} is used in operations that must succeed or fail as a whole. The application uses transactions in two critical use cases.

\subsubsection{Sending a Message}

When a user sends a message, a notification must also be created for the recipient. Both records are written inside a single transaction. If either write fails, neither is committed, preventing a state where a message exists without its notification or vice versa.

\begin{lstlisting}[style=code, language={[Sharp]C}, caption=CreateMessageUseCase — transactional write]
const savedMessage = await this.dataSource.transaction(async (manager) => {
  const message = new MessageModel();
  message.id = v4();
  message.content = dto.content;
  message.sender = sender;
  message.recipient = recipient;

  const createdMessage = await manager.save(message);

  const notification = new NotificationModel();
  notification.title = `New message from ${sender.firstname} ${sender.surname}`;
  notification.message = dto.content.substring(0, 100);
  notification.user = recipient;

  await manager.save(notification);

  return createdMessage;
});
\end{lstlisting}

\subsubsection{Deleting a User}

Deleting a user requires removing all of their related data (attachments, favourites, messages, notifications, listings) in the correct order due to foreign key constraints. All deletes run inside a single transaction so that a partial failure cannot leave orphaned records in the database.

\begin{lstlisting}[style=code, language={[Sharp]C}, caption=DeleteUserUseCase — cascading transactional delete]
await this.dataSource.transaction(async (manager) => {
  // delete attachments belonging to the user's listings
  await manager.delete(AttachmentModel, { listing: { id: In(listingIds) } });
  await manager.delete(FavoriteModel, { listing: { id: In(listingIds) } });

  // delete user-level relations
  await manager.delete(FavoriteModel, { user: { id: dto.id } });
  await manager.delete(MessageModel, { sender: { id: dto.id } });
  await manager.delete(MessageModel, { recipient: { id: dto.id } });
  await manager.delete(NotificationModel, { user: { id: dto.id } });
  await manager.delete(ListingModel, { user: { id: dto.id } });
  await manager.delete(UserModel, { id: dto.id });
});
\end{lstlisting}

% ─── 5. Logging ──────────────────────────────────────────────────────────────
\section{Logging}

\subsection{Overview}

The application uses Winston as its logging library, configured with two transports: a formatted console output for local development and a Loki transport that ships logs to Grafana Loki in production.

\subsection{Log Levels}

Winston is configured to use \texttt{info} level in production and \texttt{debug} in development. The following levels are used throughout the application:

\begin{table}[H]
    \centering
    \begin{tabular}{@{}lll@{}}
        \toprule
        \textbf{Level} & \textbf{Method} & \textbf{When used} \\
        \midrule
        error & \texttt{logger.error()} & Uncaught exceptions, HTTP errors (caught by ExceptionFilter) \\
        warn  & \texttt{logger.warn()}  & Recoverable issues \\
        info  & \texttt{logger.log()}   & Startup events, successful operations \\
        debug & \texttt{logger.debug()} & Detailed flow information (development only) \\
        verbose & \texttt{logger.verbose()} & Fine-grained tracing \\
        \bottomrule
    \end{tabular}
    \caption{Log levels and their usage}
\end{table}

\subsection{Exception Filter}

All unhandled exceptions are caught by the global \texttt{AllExceptionsFilter}. It logs the HTTP method, URL, response status code, error message, and full stack trace at \texttt{error} level before returning a structured JSON response to the client.

\begin{lstlisting}[style=code, language={[Sharp]C}, caption=AllExceptionsFilter logging an exception]
this.logger.error(
  `${request.method} ${request.url} → ${status}: ${message}`,
  stack,
  'ExceptionFilter',
);
\end{lstlisting}

\subsection{Loki Transport Configuration}

The Loki transport is only activated when the \texttt{LOKI\_URL} environment variable is set. All logs are labelled with \texttt{app: "trade2-nestjs"} so they can be filtered independently of other applications running on the same Loki instance.

\begin{lstlisting}[style=code, language={[Sharp]C}, caption=Loki transport setup in AppLogger]
const lokiTransport = new LokiTransport({
  host: lokiUrl,
  labels: { app: 'trade2-nestjs' },
  json: true,
  format: winston.format.json(),
  replaceTimestamp: true,
  onConnectionError: (err) => console.error('Loki connection error:', err),
});
\end{lstlisting}

\subsection{Loki Configuration}

Loki runs in monolithic mode (\texttt{-target=all}) with filesystem storage. The Write-Ahead Log (WAL) is enabled so that in-flight log entries are not lost if the container crashes. The container runs as root (\texttt{user: "0"}) to avoid permission errors on the Docker volume, and the image version is pinned to \texttt{3.4.2} to prevent silent breaking changes from a \texttt{:latest} update.

\begin{lstlisting}[style=code, language=yaml, caption=Loki ingester configuration]
ingester:
  wal:
    enabled: true
    dir: /loki/wal
  chunk_idle_period: 5m
  chunk_retain_period: 1m
  lifecycler:
    final_sleep: 0s
\end{lstlisting}

\subsection{Viewing Logs in Grafana}

Logs can be queried and filtered in Grafana's Explore view by selecting the Loki datasource. Useful queries:

\begin{table}[H]
    \centering
    \begin{tabular}{@{}ll@{}}
        \toprule
        \textbf{Filter} & \textbf{LogQL query} \\
        \midrule
        All logs & \texttt{\{app="trade2-nestjs"\}} \\
        Errors only & \texttt{\{app="trade2-nestjs"\} |= `error`} \\
        Specific endpoint & \texttt{\{app="trade2-nestjs"\} |= `/users`} \\
        By log level & \texttt{\{app="trade2-nestjs"\} | json | level=`error`} \\
        \bottomrule
    \end{tabular}
    \caption{Example LogQL queries in Grafana Explore}
\end{table}

Grafana's time range picker allows filtering logs to any custom window (e.g.\ last 15 minutes, last 24 hours). The log panels on the dashboard display the most recent entries in real time and can be expanded to reveal the full JSON payload of each log line.

\begin{figure}[H]
    \centering
    \fbox{\rule{0pt}{6cm}\rule{12cm}{0pt}}
    \caption{Grafana Explore view — filtering error logs with LogQL}
    \label{fig:loki-explore}
\end{figure}

% ─── 6. API Documentation ─────────────────────────────────────────────────────
\section{API Documentation}

The backend exposes a Swagger UI at \texttt{/documentation} when \texttt{SWAGGER\_ENABLED=true}. The full OpenAPI 3.0 specification is available in machine-readable form at \texttt{/documentation-json} and \texttt{/documentation-yaml}, making it importable directly into tools such as Postman or Insomnia.

\begin{figure}[H]
    \centering
    \fbox{\rule{0pt}{6cm}\rule{12cm}{0pt}}
    \caption{Swagger UI showing available endpoints}
    \label{fig:swagger}
\end{figure}

% ─── 7. Conclusion ────────────────────────────────────────────────────────────
\section{Conclusion}

\begin{table}[H]
    \centering
    \begin{tabular}{@{}p{7cm}p{7cm}@{}}
        \toprule
        \textbf{Requirement} & \textbf{Implementation} \\
        \midrule
        Zero-downtime updates & Dokploy rolling updates triggered by GitHub Actions after image push \\
        Monitoring with alerts & Prometheus scraping every 5s, Grafana dashboard with 10 panels \\
        Data consistency & TypeORM transactions on message creation and user deletion \\
        Multi-level logging & Winston with error / warn / info / debug / verbose levels \\
        Log visualisation \& filtering & Grafana Loki with LogQL filtering by level, endpoint, time range \\
        Technical documentation & This document and the backend README \\
        \bottomrule
    \end{tabular}
    \caption{Requirements coverage summary}
\end{table}

\end{document}
